% 需要的宏包（如导言区已包含，可不重复添加）：
% \usepackage{graphicx}
% \usepackage{booktabs}
% \usepackage{multirow}
% \usepackage{makecell}
% \usepackage{float}
% \usepackage{subcaption}
% \usepackage[table]{xcolor}
% \definecolor{bestcolor}{HTML}{F0D0D0}
% \definecolor{secondcolor}{HTML}{D0E8F0}

\section{实验设计与结果分析}
\label{sec:experiments}

为系统验证所提 FL-DQN 任务卸载方法的有效性，本节按照“预测层验证—端到端性能验证—鲁棒性与可解释性分析—模块贡献分析”的逻辑展开实验。首先，通过并发开销预测精度、训练收敛过程、置信校准结果和异构环境下的预测误差，验证聚类式联邦学习能否为节点侧决策提供可靠的并发风险先验；其次，通过平均任务时延、系统吞吐量、SLA 完成率和 P99 尾时延等指标，分析预测先验是否能够有效转化为实际卸载收益；随后，从链上摘要陈旧性、动作选择偏好和效率--公平性折中三个角度进一步评估方法的鲁棒性与可解释性；最后，通过消融实验分析联邦预测模块、并发惩罚项、经验回放机制和奖励塑形项对整体性能的具体贡献。

需要说明的是，受限于真实在轨天基边缘计算环境难以直接部署和复现实验，本文实验部分采用模型驱动的离散事件仿真方式进行验证。仿真参数依据前文建立的任务层时延模型、能耗模型和并发开销模型设定，并结合图像处理类任务与数值计算类任务构造不同并发强度下的资源竞争场景。实验结果主要用于比较不同策略在相同约束条件下的相对性能差异，而不代表某一具体星载平台的绝对实测性能。

\subsection{实验场景、对比方法与评价指标}
\label{subsec:exp_setup}

实验场景模拟由若干卫星终端 ST、候选卫星边缘节点 SEN、联盟链控制平面和联邦协调单元共同组成的天基分布式计算系统。终端侧任务按照离散事件方式持续到达，每个任务包含计算量、输入数据大小、时延约束、任务类型和优先级等属性；候选边缘节点具有不同的计算能力、队列长度、背景负载和并发任务数；链上公共摘要以固定周期更新，用于提供候选节点的平均负载、负载方差、信誉评分和模型版本信息。为模拟天基网络中的状态滞后现象，实验进一步设置可控的链上摘要延时，并考察不同延时条件下各策略的性能变化。

并发开销预测数据由两类代理任务构造：一类为图像处理任务，用于模拟遥感预处理、轻量识别和视觉感知类边缘负载；另一类为数值计算任务，用于模拟输入规模较小但计算强度较高的计算密集型负载。每条样本记录任务类型、输入规模、估计计算量、节点 CPU 状态、并发任务数、执行中任务剩余工作量、上下文切换强度以及实际完成时延，并根据前文定义的并发开销表达式构造监督标签。这样既能够覆盖数据搬移和计算争用同时存在的混合场景，也能够避免训练样本过度偏向单一任务类型。

对比方法包括以下几类：\textbf{FL-DQN} 表示本文提出的完整方法；\textbf{DQN-only} 表示不引入联邦预测先验，仅依靠即时状态训练 DQN 策略；\textbf{FL-only} 表示使用联邦预测结果辅助启发式卸载，但不进行强化学习策略优化；\textbf{Centralized-DQN} 表示在理想条件下能够获得全局实时状态的集中式 DQN，可视为性能上界参考；\textbf{Heuristic} 表示基于当前负载、链路质量和队列长度的启发式卸载策略。对于预测层实验，额外比较 \textbf{Clustered FL}、\textbf{FedAvg}、\textbf{Centralized} 和 \textbf{Local-only} 四类训练方式，以验证聚类式联邦学习在 Non-IID 数据分布下的优势。

评价指标分为预测层指标和决策层指标。预测层采用 MAE、RMSE 和预测--真实拟合关系衡量并发开销预测精度，同时通过置信校准结果分析模型输出不确定性的有效性。决策层采用平均任务时延、系统吞吐量、SLA 完成率、P99 尾时延、在线决策开销、跨轮次波动和 Jain 公平性指数等指标衡量卸载策略的综合表现。其中，平均任务时延反映整体服务效率，P99 尾时延反映极端拥塞条件下的稳定性，SLA 完成率反映任务时延约束满足能力，Jain 公平性指数用于衡量任务分配在不同候选节点之间的均衡程度。

\subsection{并发开销预测模块性能分析}
\label{subsec:pred_results}

节点侧卸载策略能否提前规避潜在拥塞，很大程度上取决于并发开销预测模块是否能够提供稳定可靠的风险先验。因此，本节首先对预测模块进行独立验证。图~\ref{fig:pred_perf_enhanced} 从联邦训练收敛过程、预测--真实散点密度、置信度校准结果以及不同异构强度下的测试误差四个角度展示了并发开销预测模块的综合性能。

\begin{figure}[H]
    \centering
    \includegraphics[width=0.97\textwidth]{images/4-new-6.png}
    \caption{并发开销预测模块的多维性能结果}
    \label{fig:pred_perf_enhanced}
\end{figure}

由图~\ref{fig:pred_perf_enhanced}(a) 可见，随着通信轮次增加，各类训练方式的验证集 MAE 均逐步下降，说明模型能够从历史执行日志中学习并发开销与任务特征、节点状态和资源竞争程度之间的映射关系。其中，Clustered FL 的误差下降速度明显快于 FedAvg 和 Local-only，并且在训练后期达到更低的稳定误差水平。这表明，在节点负载模式和任务分布存在差异的情况下，先根据负载模式进行聚类，再在簇内执行模型聚合，能够有效缓解统一全局模型在多峰分布下的偏移问题。

图~\ref{fig:pred_perf_enhanced}(b) 给出了预测并发开销与真实并发开销之间的散点密度关系。大多数样本集中分布在对角线附近，说明预测值与真实值总体保持较高一致性。当真实并发开销较低时，样本分布相对紧凑，预测误差较小；当真实并发开销增大时，散点分布略有扩散，说明高并发条件下资源争用的不确定性更强，但整体仍能保持较好的趋势跟踪能力。该结果说明，预测模型不仅能够学习平均意义上的并发开销变化，也能够对不同负载区间下的风险差异做出有效区分。

图~\ref{fig:pred_perf_enhanced}(c) 展示了预测不确定性与实际绝对误差之间的关系。可以看到，经验误差总体随预测不确定性增加而上升，说明模型输出的置信信息具有一定校准意义。换言之，当预测模型对某些样本给出较高不确定性时，这些样本通常也更容易产生较大的实际误差。该结果对于后续 DQN 决策具有重要意义：DQN 不仅可以利用并发开销点估计值，还可以利用不确定性信息判断当前预测是否可靠，从而避免在高风险、高不确定性状态下过度依赖单一预测值。

图~\ref{fig:pred_perf_enhanced}(d) 进一步比较了不同数据异构强度下的测试集 MAE。可以看出，在 IID 场景下，各方法之间的差距相对有限；但随着数据分布从轻度 Non-IID 增强到重度 Non-IID，FedAvg 和 Local-only 的误差增长更为明显，而 Clustered FL 的误差上升幅度较小，表现出更好的异构鲁棒性。该现象说明，标准 FedAvg 在节点负载模式差异较大时容易受到全局平均效应影响，而聚类式联邦学习能够通过簇内聚合保留局部负载特征，从而为节点侧决策提供更加稳定的并发风险估计。

综合图~\ref{fig:pred_perf_enhanced} 的结果可以看出，聚类式联邦学习不仅提高了并发开销预测精度，也提升了模型在异构环境下的收敛稳定性和置信表达能力。这为后续 FL-DQN 将预测先验引入状态空间和奖励函数提供了基础支撑。

\subsection{不同并发强度下的端到端卸载性能分析}
\label{subsec:e2e_results}

在验证预测模块有效之后，进一步需要回答的问题是：预测得到的并发风险先验能否真正转化为任务级卸载性能提升。为此，本文在不同平均并发任务数下，对 FL-DQN 与多种基线方法在平均任务时延、系统吞吐量、SLA 完成率和 P99 尾时延四项指标上的表现进行对比，结果如图~\ref{fig:service_quality} 所示。

\begin{figure}[H]
    \centering
    \includegraphics[width=0.97\textwidth]{images/4-new-7.png}
    \caption{不同并发强度下的卸载决策性能与稳定性}
    \label{fig:service_quality}
\end{figure}

由图~\ref{fig:service_quality}(a) 可见，随着平均并发任务数增加，所有方法的平均任务时延均呈上升趋势。这是因为并发任务增多后，候选节点内部的 CPU 时间片争用、缓存冲突、内存带宽竞争和 I/O 排队现象更加明显，导致任务端到端完成时间持续增加。然而，不同方法的时延增长斜率存在明显差异。Heuristic 和 DQN-only 在高并发区间时延上升较快，说明仅依赖当前负载或即时观测状态难以准确判断候选节点的未来拥塞风险。相比之下，FL-DQN 的时延增长更为平缓，说明引入并发开销预测先验后，策略能够在任务卸载前提前识别潜在拥塞节点，从而减少将任务发送到“表面空闲但实际竞争严重”节点的概率。

图~\ref{fig:service_quality}(b) 给出了系统吞吐量随并发强度变化的情况。随着并发任务数增加，各方法吞吐量均出现不同程度下降，这是由于节点队列积压和任务执行时间延长降低了单位时间内的任务完成数量。Centralized-DQN 由于假设能够获得全局实时状态，因此整体吞吐量最高，可视为理想上界。FL-DQN 虽略低于 Centralized-DQN，但明显优于 DQN-only、FL-only 和 Heuristic，说明在不依赖全局实时透明信息的前提下，所提方法通过“链上摘要 + 本地观测 + 联邦预测先验”的组合，已经能够获得接近集中式策略的任务处理能力。

图~\ref{fig:service_quality}(c) 展示了 SLA 完成率随并发强度变化的趋势。在低并发条件下，各方法之间的差距相对较小，说明系统资源较充足时，不同策略都较容易满足任务时延约束。但当平均并发任务数持续增加后，Heuristic 和 DQN-only 的完成率下降更快，表明其在高负载条件下更容易产生拥塞误判和任务超时。FL-DQN 的 SLA 完成率下降最慢，说明其能够通过预测先验提前规避高风险卸载动作，从而在资源竞争增强时仍保持较高的任务完成质量。

图~\ref{fig:service_quality}(d) 从尾部性能角度进一步验证了上述结论。P99 尾时延反映了最差一部分任务的完成情况，对时延敏感型天基业务尤为重要。随着并发任务数增加，各方法的 P99 尾时延均明显上升，但 FL-DQN 的增长幅度相对较小。相比 DQN-only 和 Heuristic，FL-DQN 在高并发区间能够显著降低尾部时延，说明其不仅改善了平均意义上的服务质量，也有效抑制了极端拥塞样本对系统稳定性的影响。

总体来看，图~\ref{fig:service_quality} 表明，FL-DQN 的优势在中高并发场景中更加明显。这与本文方法设计目标一致：当系统处于轻载状态时，候选节点之间的资源竞争尚不显著，预测模块带来的收益有限；而当系统逐步逼近拥塞边界时，并发开销预测能够为 DQN 提供更有价值的前瞻信息，使策略在平均时延、吞吐量、SLA 完成率和尾部时延方面均表现出更好的稳定性。

\subsection{鲁棒性、动作偏好与多目标折中分析}
\label{subsec:robustness_results}

仅从平均时延或完成率角度评价卸载策略，仍不足以全面说明方法的可部署性。天基网络中链路窗口有限、链上摘要存在确认延迟，节点状态不可避免地具有陈旧性；同时，卸载策略还应具备可解释的动作选择行为，并在效率、能耗和节点均衡之间取得合理折中。基于此，本文进一步从链上陈旧性敏感性、动作选择偏好和效率--公平性折中三个角度进行分析，结果如图~\ref{fig:robustness_tradeoff} 所示。

\begin{figure}[H]
    \centering
    \includegraphics[width=0.96\textwidth]{images/4-new-8.png}
    \caption{鲁棒性、动作偏好与多目标折中结果}
    \label{fig:robustness_tradeoff}
\end{figure}

图~\ref{fig:robustness_tradeoff}(a) 展示了链上摘要延时增加时各方法的平均任务时延变化。可以看到，随着链上摘要延时从较低水平逐步增加，各方法的平均时延均有所上升。这是因为候选节点公共摘要越陈旧，终端侧获得的节点负载状态越可能偏离真实运行状态，从而增加错误卸载的概率。然而，FL-DQN 的退化幅度明显小于 DQN-only、FL-only 和 Heuristic。该结果说明，FL-DQN 并不完全依赖链上摘要本身，而是能够通过本地实时观测和联邦预测模型对状态陈旧性进行一定补偿。当链上信息不能准确反映当前节点状态时，预测模块仍可根据历史负载模式和并发环境特征估计潜在风险，从而降低摘要延迟对决策质量的影响。

图~\ref{fig:robustness_tradeoff}(b) 给出了 FL-DQN 在不同候选节点负载风险区间下的动作选择热力图。该热力图通过统计策略收敛后在测试集中的动作选择频次，并对同一风险区间内的动作次数进行归一化得到。可以看出，在负载风险较低时，策略更倾向于选择本地执行或风险较低的候选边缘节点；随着候选节点负载风险逐步升高，动作选择概率逐渐向其他节点分散，表明智能体会根据风险变化动态调整卸载目标，而不是长期固定选择某一个高性能节点。该结果说明，FL-DQN 学到的并不是简单的“选择当前负载最低节点”规则，而是一种结合任务属性、链上摘要和并发风险预测的自适应卸载行为。

图~\ref{fig:robustness_tradeoff}(c) 进一步展示了不同方法在能时积 EDP 与 Jain 公平性指数之间的折中关系，其中气泡大小表示对应方法的任务完成率。EDP 越低，说明方法在时延和能耗之间具有更好的综合效率；Jain 公平性指数越高，说明任务在不同节点之间分配越均衡。可以看到，Centralized-DQN 由于拥有理想全局状态，在公平性和效率方面表现较好，但其依赖强全局实时信息假设，实际部署难度较高。FL-DQN 在 EDP 和公平性之间取得了较优折中，虽然其公平性略低于集中式上界，但明显优于 DQN-only、FL-only 和 Heuristic，并且仍保持较高任务完成率。相比之下，Heuristic 和 DQN-only 更容易将任务集中到少数表面负载较低的节点，从而导致公平性下降和能时积增大。

综合图~\ref{fig:robustness_tradeoff} 可以看出，FL-DQN 的优势不仅体现在平均性能指标上，还体现在对链上状态滞后的适应能力、动作选择行为的合理性以及效率--公平性多目标折中能力上。这说明所提方法更符合天基分布式计算环境中“信息不完全、状态有时滞、节点负载动态变化”的实际部署需求。

\subsection{消融实验与模块贡献分析}
\label{subsec:ablation_results}

为了进一步分析完整方法中各组成模块的具体作用，本文设计消融实验，分别移除联邦预测模块、并发惩罚项、经验回放机制和奖励塑形项，并与完整 FL-DQN 进行对比。除被移除模块外，其余训练参数、任务到达过程、候选节点配置和评价指标均保持一致，以保证不同消融版本之间具有可比性。为便于直接对比，表~\ref{tab:ablation_summary} 汇总了平均并发任务数为 $16$ 时的关键指标，其中 Avg. Delay 和 SLA 分别来自图~\ref{fig:ablation_results}(a)(b) 的高并发点，Online Cost 和 Std. Dev. 分别对应图~\ref{fig:ablation_results}(c)(d) 的柱状统计结果。表中红色表示最优结果，蓝色表示次优结果。

\begin{table}[H]
\centering
\footnotesize
\caption{FL-DQN 各模块消融实验结果}
\label{tab:ablation_summary}
\setlength{\tabcolsep}{2.3mm}
\renewcommand{\arraystretch}{1.18}
\begin{tabular}{lccccp{3.1cm}}
\toprule
\rowcolor{gray!10}
方法变体
& Avg. Delay $\downarrow$
& SLA $\uparrow$
& Online Cost $\downarrow$
& Std. Dev. $\downarrow$
& 结果说明 \\
\midrule
FL-DQN
& \cellcolor{bestcolor}\textbf{209}
& \cellcolor{bestcolor}\textbf{91.2}
& 3.8
& \cellcolor{bestcolor}\textbf{5.1}
& 完整方法，综合性能最优 \\

-FL 预测
& 284
& 80.1
& \cellcolor{secondcolor}\textbf{3.1}
& 9.6
& 缺少并发风险先验，退化最明显 \\

-并发惩罚
& 243
& 84.7
& 3.6
& 8.1
& 拥塞规避约束减弱 \\

-经验回放
& 256
& 86.4
& \cellcolor{bestcolor}\textbf{2.4}
& 11.2
& 开销最低但训练波动最大 \\

-奖励塑形
& \cellcolor{secondcolor}\textbf{238}
& \cellcolor{secondcolor}\textbf{87.9}
& 3.4
& \cellcolor{secondcolor}\textbf{7.3}
& 收敛引导能力减弱 \\
\bottomrule
\end{tabular}
\end{table}

进一步地，图~\ref{fig:ablation_results} 从平均时延、SLA 完成率、在线决策开销和跨轮次波动四个角度展示了消融实验结果。为了便于排版，图~\ref{fig:ablation_results} 由四个 Python 绘制的子图组合而成。

\begin{figure}[H]
    \centering
    \begin{subfigure}[t]{0.48\textwidth}
        \centering
        \includegraphics[width=\linewidth]{figures/fig4_9_a_avg_delay_cn.png}
        \caption{平均时延随并发度变化}
        \label{fig:ablation_delay}
    \end{subfigure}
    \hfill
    \begin{subfigure}[t]{0.48\textwidth}
        \centering
        \includegraphics[width=\linewidth]{figures/fig4_9_b_sla_cn.png}
        \caption{SLA 完成率随并发度变化}
        \label{fig:ablation_sla}
    \end{subfigure}

    \vspace{0.25cm}
    \begin{subfigure}[t]{0.48\textwidth}
        \centering
        \includegraphics[width=\linewidth]{figures/fig4_9_c_overhead_cn.png}
        \caption{在线决策开销}
        \label{fig:ablation_overhead}
    \end{subfigure}
    \hfill
    \begin{subfigure}[t]{0.48\textwidth}
        \centering
        \includegraphics[width=\linewidth]{figures/fig4_9_d_variance_cn.png}
        \caption{跨轮次波动}
        \label{fig:ablation_variance}
    \end{subfigure}
    \caption{消融实验结果}
    \label{fig:ablation_results}
\end{figure}

图~\ref{fig:ablation_results}(a) 展示了不同消融版本在平均并发任务数增加时的平均时延变化。可以看到，完整 FL-DQN 在所有并发水平下均保持最低平均时延；去除 FL 预测模块后，平均时延显著上升，且并发任务数越高，性能劣化越明显。这说明联邦预测模块是完整方法取得性能优势的核心来源。没有预测先验时，DQN 只能依赖当前观测状态和链上摘要进行反应式决策，难以及时识别候选节点的潜在并发风险，因此在高并发条件下更容易发生错误卸载。相比之下，去除并发惩罚项后时延也有所增加，但劣化程度低于去除 FL 预测模块，说明并发惩罚主要起到强化拥塞规避偏好的作用，而预测模块则决定了策略是否具备风险前瞻能力。

图~\ref{fig:ablation_results}(b) 给出了不同消融版本的 SLA 完成率变化趋势。完整 FL-DQN 的完成率始终最高，并且随着并发增强下降较为平缓；去除 FL 预测模块后，SLA 完成率下降最明显，说明缺少并发风险估计会直接影响任务是否能够在时延阈值内完成。去除并发惩罚项后，SLA 完成率也明显低于完整方法，表明仅将预测值作为状态输入仍不足以完全引导策略规避拥塞，还需要在奖励函数中显式体现并发开销代价，使智能体在训练过程中形成对高风险动作的惩罚记忆。

图~\ref{fig:ablation_results}(c) 比较了各版本的在线决策开销。完整 FL-DQN 的在线开销略高于去除 FL 预测模块或去除经验回放的简化版本，这是因为完整方法需要在决策前调用预测模型，并进行 DQN 动作价值计算。然而，其决策开销仍维持在毫秒级范围内，满足节点侧实时决策需求。值得注意的是，去除经验回放后虽然在线开销最低，但这并不意味着整体效果更优，因为经验回放主要影响训练稳定性和样本利用效率，其贡献更多体现在长期性能和跨轮次波动上。

图~\ref{fig:ablation_results}(d) 展示了不同版本的跨轮次波动情况。完整 FL-DQN 的标准差最低，说明其在多次独立运行中表现更稳定；去除经验回放后波动最大，表明直接使用连续交互样本更新网络会增强训练样本相关性，导致策略更容易震荡；去除 FL 预测模块后波动也明显增大，说明缺少稳定风险先验会使策略在相似观测状态下产生更不稳定的动作选择。去除奖励塑形项后波动介于两者之间，说明合理的奖励设计有助于提高策略收敛质量，但其影响弱于 FL 预测和经验回放。

综合消融实验结果可以看出，FL 预测模块决定了策略的风险前瞻能力，是降低高并发时延和维持 SLA 完成率的关键；并发惩罚项增强了策略对拥塞节点的规避倾向，使预测信息能够更直接地转化为优化目标；经验回放机制主要提升训练稳定性，降低跨轮次波动；奖励塑形项则进一步改善策略收敛方向和任务完成质量。因此，完整 FL-DQN 的性能提升并非来自单一模块堆叠，而是来自预测先验、奖励约束和强化学习稳定机制之间的协同作用。

\subsection{本节小结}
\label{subsec:exp_summary}

本节围绕并发感知任务卸载方法的有效性进行了分层实验验证。预测层结果表明，聚类式联邦学习能够在不集中上传原始日志的前提下，有效学习候选节点的并发开销变化规律，并在 Non-IID 环境下保持较好的预测精度和收敛稳定性。端到端卸载结果表明，将并发开销预测值及其置信信息引入 DQN 后，节点侧策略在平均任务时延、系统吞吐量、SLA 完成率和 P99 尾时延等指标上均优于分布式基线方法，尤其在中高并发区间优势更加明显。鲁棒性分析进一步说明，FL-DQN 能够在链上摘要存在陈旧性的情况下维持较平缓的性能退化，并形成符合风险变化规律的动作选择行为。消融实验则验证了联邦预测模块、并发惩罚项、经验回放机制和奖励塑形项对整体性能提升的不同贡献。

总体而言，实验结果说明，本文所提出的 FL-DQN 方法能够在隐私受限、状态滞后和负载异构的天基分布式计算环境中，将联邦学习得到的并发风险预测有效转化为节点侧实时卸载决策收益，从而在时延、吞吐量、完成率、尾部稳定性和负载均衡之间取得更合理的综合折中。
